\chapter{实分析：从全等图形到测度与连续谱}
\label{chap:real-analysis}

实分析可以把「图形有多大」推广为勒贝格测度，也可以在 \(L^2\) 中研究函数在各点的取值。测度与距离是两种不同的信息。若 \(Q^TQ=I\)，则 \((\det Q)^2=1\)，由线性变量代换公式，对可测集 \(E\subseteq\mathbb R^n\) 有
\[
 m(QE)=\int_E |\det Q|\,dx=m(E).
\]
平移同样保测。因而欧氏全等推出相同的测度；面积相等本身不足以推出全等。

\section{保面积与保距离}

\begin{example}{相同的面积能否推出全等}
取矩形 \(R=[0,2]\times[0,1]\)。分别作 \(Q=\frac1{\sqrt5}\begin{pmatrix}2&-1\\1&2\end{pmatrix}\) 与剪切 \(S=\begin{pmatrix}1&2\\0&1\end{pmatrix}\)。比较 \(QR\)、\(SR\) 与 \(R\) 的面积、全等关系。
\end{example}
\begin{solution}
两矩阵的行列式都等于 \(1\)，因此变量代换给出
\[
 m(QR)=m(SR)=m(R)=2.
\]
由于 \(Q^TQ=I\)，它使任意两点的距离保持不变，又可逆，所以 \(QR\) 与 \(R\) 全等。剪切矩阵把 \(R\) 的两个顶点 \((0,0),(2,1)\) 分别送到 \((0,0),(4,1)\)。在 \(R\) 中，两点横坐标之差的绝对值至多 \(2\)，纵坐标之差的绝对值至多 \(1\)，因此任意两点距离都不超过 \(\sqrt5\)。然而 \(SR\) 里已有相距 \(\sqrt{17}\) 的两点，故不存在把 \(R\) 保距地映到整个 \(SR\) 的映射。这不仅说明 \(S\) 本身不保距，还排除了用别的等距变换使两个图形全等的可能。
\end{solution}

\begin{example}{椭圆区域和两条等值线之间的面积}
令 \(q_A(x,y)=5x^2+4xy+2y^2\)，\(E_c=\{(x,y):q_A(x,y)\leq c\}\)。求 \(c>0\) 时的面积，以及 \(0<a<b\) 时区域 \(\{a<q_A\leq b\}\) 的面积。说明等值线 \(q_A=c\) 的平面测度为零。
\end{example}
\begin{solution}
作 \((x,y)^T=Q(u,v)^T\)，则 \(q_A=6u^2+v^2\)，且 \(Q\) 不改变面积。再令
\[
 u=\sqrt{c/6}\,s,\qquad v=\sqrt c\,t.
\]
此时 \(6u^2+v^2\leq c\) 等价于 \(s^2+t^2\leq1\)。第二次变换的行列式为 \(c/\sqrt6\)，所以
\[
 m(E_c)=\frac{\pi c}{\sqrt6},\qquad
 m(\{a<q_A\leq b\})=m(E_b)-m(E_a)=\frac{\pi(b-a)}{\sqrt6}.
\]
这里 \(E_a\subseteq E_b\)，故可直接相减。对 \(0<h<c\)，等值线 \(q_A=c\) 包含在 \(\{c-h<q_A\leq c\}\) 中，其面积至多 \(\pi h/\sqrt6\)；让 \(h\) 趋零即得结论。第一步的 \(Q\) 是全等变换；第二步是用来算面积的伸缩，不保持全等。
\end{solution}

\section{把积分与函数的取值分布联系起来}

\begin{example}{利用等值区域计算二维高斯积分}
求 \(\int_{\mathbb R^2}e^{-q_A(x,y)}\,dx\,dy\)。若 \(c>0\)，在 \(E_c\) 中按面积均匀取一点 \(X\)，求 \(q_A(X)/c\) 的分布。
\end{example}
\begin{solution}
对每个点，\(q_A\geq0\)，而 \(e^{-q_A}=\int_{q_A}^\infty e^{-s}\,ds\)。把积分下限改成 \(0\) 并乘上条件 \(q_A\leq s\) 的示性函数，再用非负函数的 Tonelli 定理交换积分次序：
\[
 \int_{\mathbb R^2}e^{-q_A}\,dx\,dy
 =\int_0^\infty e^{-s}m(E_s)\,ds
 =\frac{\pi}{\sqrt6}\int_0^\infty se^{-s}\,ds
 =\frac{\pi}{\sqrt6}.
\]
最后一步由分部积分得到 \(\int_0^\infty se^{-s}\,ds=1\)。均匀取点时，对 \(0\leq r\leq1\)，有
\[
 \mathbb P(q_A(X)/c\leq r)=\frac{m(E_{rc})}{m(E_c)}=r.
\]
\(r=0\) 时只有原点，面积为零；\(r=1\) 时整个区域都满足条件。因此 \(q_A(X)/c\) 在 \([0,1]\) 上均匀分布。计算依靠的是正交变换下的面积不变和面积随 \(c\) 的线性增长，而非点的角度均匀。
\end{solution}

\begin{example}{旋转函数图样，什么数据随之保持}
定义 \(Uf(z)=f(Qz)\)，其中 \(f\in L^2(\mathbb R^2)\)。证明 \(U\) 是满射等距。对 \(f=2\mathbf 1_{E_1}-\mathbf 1_{E_2\setminus E_1}\)，求 \(\|Uf\|_2^2\) 和 \(m(\{z:|Uf(z)|>t\})\)（\(t\geq0\)）。相同的取值分布能否推出两个图样旋转或平移后重合？
\end{example}
\begin{solution}
变量代换 \(x=Qz\) 给出
\[
 \|Uf\|_2^2=\int_{\mathbb R^2}|f(Qz)|^2\,dz
 =\int_{\mathbb R^2}|f(x)|^2\,dx.
\]
若两个函数只在零测集上不同，则拉回后也只在零测集上不同，故 \(U\) 在 \(L^2\) 的等价类上定义良好。\(U^{-1}g(x)=g(Q^Tx)\) 给出满射性。并且
\[
 \{z:|Uf(z)|>t\}=Q^T\{x:|f(x)|>t\};
\]
两集的测度相等。由于 \(m(E_1)=m(E_2\setminus E_1)=\pi/\sqrt6\)，可得
\[
 \|Uf\|_2^2=\frac{5\pi}{\sqrt6},\qquad
 m(\{|Uf|>t\})=
 \begin{cases}
 2\pi/\sqrt6,&0\leq t<1,\\
 \pi/\sqrt6,&1\leq t<2,\\
 0,&t\geq2.
 \end{cases}
\]
反推不成立。例如 \(R=[0,2]\times[0,1]\) 与
\(C=([0,1]\cup[3,4])\times[0,1]\) 面积相同，示性函数取值都只有 \(0,1\)，相应正阈值的超水平集测度相同。但 \(R\) 的直径是 \(\sqrt5\)，\(C\) 的直径是 \(\sqrt{17}\)；等距变换保直径，所以无法使两个图样重合。即使允许相差零测集也不行：任一顶点的任意邻域内都有正面积的矩形部分，删掉零测集仍可任意逼近该顶点，因此两集的直径各自保持不变。
\end{solution}

\section{乘法算子：谱可能是一整个区间}

在 \(L^2[0,1]\) 中定义 \(Mf(t)=tf(t)\)。它像对角矩阵一样「每个位置单独乘一个数」，但在函数空间里，单点的测度为零，不能把每个位置当成一个非零坐标向量。此时谱分解应看成按取值范围切分函数。

\begin{example}{自伴算子没有非零特征向量}
证明 \(M\) 有界自伴、非紧，且谱为 \([0,1]\)。证明 \(M\) 没有特征值。
\end{example}
\begin{solution}
因为 \(0\leq t\leq1\)，有 \(\|Mf\|_2\leq\|f\|_2\)。内积中乘子 \(t\) 为实数，所以 \(\langle Mf,g\rangle=\langle f,Mg\rangle\)。若 \(Mf=\lambda f\)，则 \((t-\lambda)f(t)=0\) 几乎处处。除至多一个点以外，\(t-\lambda\) 不为零，所以 \(f=0\) 几乎处处；不存在非零特征函数。

若 \(\lambda\notin[0,1]\)，则 \(1/(t-\lambda)\) 在 \([0,1]\) 上有界，乘以它就得到 \(M-\lambda I\) 的有界逆。若 \(\lambda\in[0,1]\)，取 \(I_h=[0,1]\cap(\lambda-h,\lambda+h)\)（\(0<h<1\)），令 \(f_h=\mathbf1_{I_h}/\sqrt{m(I_h)}\)。则
\[
 \|f_h\|_2=1,\qquad \|(M-\lambda I)f_h\|_2\leq h.
\]
若有有界逆，便会推出 \(1\leq\|(M-\lambda I)^{-1}\|h\)，让 \(h\) 趋零即矛盾。因此每个 \(\lambda\in[0,1]\) 都属于谱。

再取 \([1/2,1]\) 中两两不交、长度为正的区间 \(J_n=(1/2+2^{-(n+2)},1/2+2^{-(n+1)})\)，其中 \(n\geq1\)。令 \(g_n=\mathbf1_{J_n}/\sqrt{m(J_n)}\)。各 \(g_n\) 的范数为 \(1\)，各 \(Mg_n\) 的范数至少为 \(1/2\)，且支撑两两不交。因此 \(n\ne k\) 时
\[
 \|Mg_n-Mg_k\|_2^2=\|Mg_n\|_2^2+\|Mg_k\|_2^2\geq\frac12.
\]
序列 \((Mg_n)\) 无收敛子列，故 \(M\) 非紧。此例说明无穷维有界自伴算子不必能按普通特征向量逐个对角化。
\end{solution}

\begin{example}{把一个谱区间上的分量取出来}
对上述 \(M\)，给定可测集 \(B\subseteq[0,1]\)，定义 \(P_B f=\mathbf1_B f\)。验证它是正交投影。令 \(f(t)=\sqrt3\,t\)，求 \(\|P_{[a,b]}f\|_2^2\) 与 \(\langle Mf,f\rangle\)。再证明 \(M\) 与 \(M_{1-t}:f(t)\mapsto(1-t)f(t)\) 的二次型等值集可由满射等距对应。
\end{example}
\begin{solution}
由 \(\mathbf1_B^2=\mathbf1_B\) 可知 \(P_B^2=P_B\)，并且 \(f=P_Bf+(f-P_Bf)\)，两项支撑互不重叠而正交。因此 \(P_B\) 正是到「在 \(B\) 外为零」的闭子空间的正交投影。对 \(0\leq a\leq b\leq1\)，
\begin{align*}
 \|P_{[a,b]}f\|_2^2&=\int_a^b 3t^2\,dt=b^3-a^3,\\
 \|f\|_2^2&=1,\qquad
 \langle Mf,f\rangle=\int_0^1 3t^3\,dt=\frac34.
\end{align*}
可把 \(3t^2\,dt\) 看作此单位向量在各段谱区间上的「平方长度分配」：单点分量为零，整个区间却可以有正的分量。这是连续谱中用投影表达分解的具体方式。

令 \(Vf(t)=f(1-t)\)。换元 \(s=1-t\) 给出 \(\|Vf\|_2=\|f\|_2\)，且 \(V^{-1}=V\)。逐点计算得
\[
 (V^{-1}MVf)(t)=(1-t)f(t),
 \qquad
 \langle M_{1-t}f,f\rangle=\langle MVf,Vf\rangle.
\]
所以 \(V\) 把前一二次型的任一等值集保距地送到后一二次型的相应等值集，反向用 \(V^{-1}\) 即可。整个过程不需要选出一个特征向量。全等关系来自满射等距所给出的共轭；是否有普通特征向量基是另一个问题。
\end{solution}
